% ================================================================
% ch01_biography.tex — अध्याय १: संत दरिया साहब: जीवन वृत्त आ साहित्यिक योगदान
% ================================================================

\cleardoublepage
\section*{अध्याय १: जीवन वृत्त आ साहित्यिक योगदान}
\addcontentsline{toc}{section}{अध्याय १: जीवन वृत्त आ साहित्यिक योगदान}

\subsection*{१.१ जन्म आ प्रारंभिक जीवन}
\addcontentsline{toc}{subsection}{१.१ जन्म आ प्रारंभिक जीवन}

संत दरिया साहब के जन्म 1674 ईस्वी (विक्रम संवत 1731) में बिहार के रोहतास/भोजपुर क्षेत्र के प्रसिद्ध गाँव धरकंधा {\engfont (Dharkandha)} में भइल रहे। इनकर पिता के नाम पीरन शाह आ माता के नाम टेकनी देवी रहे। दरिया साहब के परिवार दरजी व्यवसाय से जुड़ल रहे। ऊ बेहद सरल, स्वावलंबी आ निष्ठावान स्वभाव के रहलें।

\subsection*{१.२ आध्यात्मिक चेतना आ निर्गुण भक्ति}
\addcontentsline{toc}{subsection}{१.२ आध्यात्मिक चेतना आ निर्गुण भक्ति}

दरिया साहब कबीरपंथी आ संत-मत के महान परंपरा के प्रतिनिधि संत-कवि बाड़ें। ऊ बाह्य आडंबर, मूर्ति-पूजा, जाति-पाँति के भेद-भाव आ धार्मिक कट्टरता के कड़ा विरोध कइलीं। उनकर विचार रहे कि ईश्वर कवनो मंदिर या मस्जिद में ना, बलुक हर मनुष्य के हृदय में विराजमान बाड़ें।

\subsection*{१.३ रचना संसार आ प्रमुख ग्रंथ}
\addcontentsline{toc}{subsection}{१.३ रचना संसार आ प्रमुख ग्रंथ}

संत दरिया साहब लगभग 15,000 से अधिक पदों आ दोहों के रचना कइलीं। इनकर प्रमुख ग्रंथ निम्नलिखित बाड़ें:

\begin{itemize}
  \item \textbf{दरिया सागर:} इनकर सबसे विशाल आ महत्वपूर्ण दार्शनिक ग्रंथ।
  \item \textbf{ज्ञान दीपक:} ज्ञान आ भक्ति के गूढ़ रहस्य बतावे वाला ग्रंथ।
  \item \textbf{भक्ति हेतु:} भक्ति के सहज मार्ग के निरूपण।
  \item \textbf{प्रेम गान आ निरगुनी पद संग्रह:} भोजपुरी लोक-राग में रचल भक्ति पद।
\end{itemize}

\subsection*{१.४ अठारहवीं सदी के भोजपुरी भाषा के महत्व}
\addcontentsline{toc}{subsection}{१.४ अठारहवीं सदी के भोजपुरी भाषा के महत्व}

दरिया साहब के काव्य भाषा 18वीं शताब्दी के विशुद्ध आ सहज लोक-भोजपुरी ह। एह में पूरबी अवधी आ तत्कालीन लोक-बोलियों के सम्मिश्रण बा, जे एह बात के प्रमाण ह कि 18वीं सदी में भोजपुरी में कितना समृद्ध आध्यात्मिक आ दार्शनिक साहित्य रचल जा रहल रहे।
